\documentclass[11pt,a4paper]{article}
\usepackage[margin=25mm]{geometry}
\usepackage{amsmath,amssymb,amsthm,hyperref}
\setlength{\emergencystretch}{2em}
\newtheorem{theorem}{Theorem}
\newcommand{\gen}[1]{\langle#1\rangle}
\title{Global eligible pairs behind a six-point head}
\author{Human Compression project: L044 continuation}
\date{}
\begin{document}\maketitle
Retain the full standing assumptions and smaller-degree alternating
induction of the accepted critical-head and paired-head theorems.
In particular $G\le A_n$ is transitive with a connected string
C-marking, $n\ge12$, and every maximal ambient omission is
intransitive. The first three generators form the actual supported
six-point head of source Table 2(2) or (3), in an extendable orientation.
The complete exceptional-head interface is required: on
$D=\{1,\ldots,6\}$ the head is transitive;
$\tau=\rho_3=(6\ b)\delta$, $\sigma=\rho_4$ fixes $D$,
$H=G_{\ge5}$ fixes $D\cup\{b\}$, and
$B=\gen{\sigma,H}$ is transitive on $\Omega\setminus D$.
Under rank failure the accepted reduction supplies $n=2r\ge18$
and the actual natural alternating tail used in the pure-head proof.
Write $K=G_{>5}$. That proof gives
\[
           x=b^\sigma\ne b,\qquad x^\tau=x,
                       \qquad K\text{ fixes }x.        \tag{1}
\]
It uses full ambient IP, not a quotient IP.

\begin{theorem}
If every colour $j\ge5$ has at least two ambient-eligible edges
anywhere in $\Omega$, then $r\le\lfloor(n-1)/2\rfloor$.
Consequently a counterexample in this exact head configuration has a
further colour $j\ge5$ with exactly one globally ambient-eligible edge.
\end{theorem}
\begin{proof}
Suppose the bound fails. Let $f$ be the number of $H$-fixed points
and $c$ its number of nontrivial orbits. The full ambient-interval
bound gives $f+c\le10$, while $f\ge7$ and $c\ge1$.
If $c=1$, every eligible pair is already in that single nontrivial
orbit, so the complete paired-head theorem applies.

Suppose $c=2$. Then $f=7$ or $8$, and $\tau$ fixes all additional
$H$-fixed points. Its action outside the seven-point head is odd.
If it preserves both nontrivial orbits, its nontrivial actions there
are free, and it is free on at least one orbit. Any eligible edge on
a free orbit has a distinct eligible translate by $\tau$: an invariant
unordered edge would be a $\tau$-pair and hence ineligible since
$\tau\in G_j$. Thus every colour with an edge on a free orbit has
a pair there; a colour with no such edge has its two edges in the
single remaining fixed orbit. The paired-head theorem again applies.

If $\tau$ exchanges the two orbits, they have the same odd size.
Their union has even size, so $f=8$. By (1), $x$ cannot lie in either
exchanged orbit, nor in $D\cup\{b\}$. It is therefore the sole
additional $H$-fixed point. Both $H$ and $\sigma$ preserve
$\{b,x\}$, contradicting transitivity of $B$. Here $x$ is first
known to be $K$-fixed; its being $H$-fixed follows only from this
exact orbit count and $\tau x=x$.

It remains that $c=3,f=7$. The nontrivial support has odd size
$n-7$. If $\tau$ preserves all three orbits and is free on at least
two, then every colour has a pair in one orbit: any edge on a free
orbit supplies that pair, and at most one fixed orbit remains.
It cannot be free on all three or act trivially on all three.
Thus the only preserved case not already covered has an odd free
action on one orbit $O$ of size $2q$, with $q$ odd, and fixes the
other two orbits pointwise. Call their union $V$ and put $v=|V|$;
then $v$ is odd and
\[
                         r-5=q+(v-3)/2.                \tag{2}
\]

Partition all tail labels into $L$, those with any ambient-eligible
edge in $O$, and $F$, those with none. These sets are disjoint and
exhaust the tail. Choose one such edge for each label in $L$ and
project to the $q$ pairs of $\tau$ in $O$. Each is a genuine edge
between different pairs. The projected edges form a forest: a cycle
would lift to an alternate path avoiding one label, using the other
selected edges and $\tau\in G_j$ inside pairs. This contradicts
ambient eligibility of that label. Hence $|L|\le q-1$.
Every eligible edge of a colour in $F$ lies in $V$, since the head
is fixed. The selected-localization theorem, using the full tail
interval and its two actual orbits in $V$, gives
$2|F|\le v-2$, and hence $|F|\le(v-3)/2$.
Together these bounds contradict (2), including $q=1$.

The final possibility is that $\tau$ exchanges two nontrivial
orbits $O,O'$ of odd size $m\ge3$ and fixes the remaining odd orbit
$V$ of size $v\ge3$. Then
\[
                         r-5=m+(v-3)/2.                \tag{3}
\]
Let $L$ be the labels with an eligible edge in $O$; equivariance
makes these exactly the labels with an eligible edge in $O'$.
Put $F=\{5,\ldots,r-1\}\setminus L$.
One chosen edge per $L$ label is a forest on $O$, so $|L|\le m-1$.
All eligible $F$-edges lie in $V$, and the selected-localization
bound gives $|F|\le(v-1)/2$. Equality is forced in both bounds
by (3).

The $m-1$ distinguished $L$ restrictions to $O$ are independent:
if one were generated by the others, its selected edge would lie
within an ambient omission orbit. The retained Whiston equality
and complete odd-degree maximal-string theorem identify their
marked image with the ordinary transposition path of $S_m$.
For $m=3$ this is simply the pair of distinct transpositions in
$S_3$; the accepted theorem covers all odd $m\ge5$, including
its degree-five certificate. Connected noncommutations force $L$
to be consecutive in the original string.

Deleting one path label leaves its two path-component orbits.
Every $F$ restriction preserves these two sets: otherwise adjoining
it to the other path generators makes the row transitive, contrary
to the chosen eligible edge. Preservation of all path cuts forces
every $F$ label to fix the row pointwise, and likewise the exchanged
row $O'$. Hence $G_F$ is actually supported on $V$.
The complement of the consecutive $L$ interval has no one-colour
hole, because $|L|=m-1\ge2$.
Apply the second selected-localization theorem to $G_F$ and $V$:
\[
                   v-1=2|F|\le v-o_{G_F}(V).
\]
Thus $G_F$ is transitive on $V$. This conclusion concerns the
actual supported subgroup; no projected intersection property was
used to obtain it.

If the first tail label $5$ belongs to $L$, then $G_F\le K$.
It centralizes $\sigma$ and has just one nontrivial orbit, namely
$V$, so $\sigma$ preserves $V$. This contradicts transitivity of
$B$. Suppose instead $5\in F$. By (1) and the exchanged rows,
$x\in V$. The set $F$ consists of an initial interval and possibly
a terminal interval. Its terminal subgroup lies in $K$, fixes $x$,
and commutes with the initial subgroup. Because their product is
transitive on $V$, the terminal subgroup fixes every point of $V$.
It is supported there, so is trivial; nonidentity of every actual
marked generator forces that terminal interval to be empty.

Consequently $L$ is terminal. Its last label $\lambda$ is at
distance at least two from the initial $F$ interval, since
$|L|\ge2$. It centralizes $G_F$ and fixes $x$ by membership in
$K$, so fixes $V$ pointwise. But all of $L$ belongs to $K$ and
is transitive on each row $O,O'$, and $\lambda$ acts nontrivially
on each. These two rows are therefore actual $K$-orbits on which
$\lambda$ acts nontrivially, whereas it fixes every point outside
them. The permutation $\sigma$ centralizes $K$, so transports its
orbits equivariantly and commutes with $\lambda$. It must preserve
$O\cup O'$. The group $H$ also preserves this union, contradicting
transitivity of $B$ for the final time.

These cases exhaust the odd action of $\tau$ outside its head.
Every tail colour has at least one ambient-eligible edge by full
transitivity and intransitivity of its omission. The contrapositive
therefore gives the stated unique edge.
\end{proof}

\paragraph{Attribution and remaining boundary.}
This completes the global-pair version of the exceptional-head step
supporting Cameron--Fernandes--Leemans--Mixer,
\emph{Highest rank of a polytope for $A_n$}, arXiv:1605.09173v1,
Proposition 5.31, physical pages 42--44. The source's orbit-deficit
strategy and all accepted head, normal-form, pure-group and finite
maximal-string inputs remain paid. The new selected-localization
bounds and equality routing replace an assumption that both edges
were already in one orbit. The fixed-point distinction in the
exchanged two-orbit case is explicit. The remaining unique-crossing
route, stronger source equality diagrams, tables and whole L044
recovery are not claimed. No new census, formal declaration,
original credit or question solution is asserted.
\end{document}
